%! Quadratic Functions Revisited — Unit 2, Lesson 2.1 — Guided Notes
\documentclass[10pt]{article}
\usepackage{precalculus-article}
\usepackage{precalculus-boxes}
\newcommand{\TallMath}[1]{$\displaystyle #1\rule[-1.4em]{0pt}{3.2em}$}

\begin{document}

\pageheader{Unit 2, Lesson 2.1}{Guided Notes}

\vspace{0.12in}

\begin{objectivebox}
By the end of this lesson, I will be able to\ldots
\begin{enumerate}[itemsep=2pt]
  \item Read the vertex, axis of symmetry, and zeros of a quadratic from its
        \blank{2cm}, \blank{1.8cm}, and \blank{2cm} forms.
  \item Compute the \blank{4.2cm} of a quadratic over an interval and
        recognize it as the slope of a \blank{2cm} line.
  \item Use \blank{1.8cm} and \blank{2cm} differences to tell a linear
        function from a quadratic one and to justify \blank{2.2cm}.
\end{enumerate}
\end{objectivebox}

\vspace{0.1in}

\boxguard
\begin{vocabbox}
Fill in each term as we define it together.
\par\vspace{2pt}
\noindent\textbf{\textcolor{plum}{Vertex form:}}\\[1pt]
\writeline\\[3pt]
\noindent\textbf{\textcolor{plum}{Average rate of change (AROC):}}\\[1pt]
\writeline\\[3pt]
\noindent\textbf{\textcolor{plum}{Secant line:}}\\[1pt]
\writeline\\[3pt]
\noindent\textbf{\textcolor{plum}{First differences:}}\\[1pt]
\writeline\\[3pt]
\noindent\textbf{\textcolor{plum}{Second differences:}}\\[1pt]
\writeline\\[3pt]
\noindent\textbf{\textcolor{plum}{Concave up / concave down:}}\\[1pt]
\writeline
\end{vocabbox}

\vspace{0.1in}

\boxguard
\begin{hookbox}
Here are the perfect squares: $0,\ 1,\ 4,\ 9,\ 16,\ 25$. The gaps between
them are $1,\ 3,\ 5,\ 7,\ 9$ --- every odd number, in order. Now find the
gaps between \textit{those} numbers. What do you get, and why is that
surprising?

\writelines{2}
\end{hookbox}

\vspace{0.1in}

\boxguard[30]
\begin{notesbox}{1. Three Forms, Three Features}
A quadratic can be written three different ways. They are the same function
--- but each one hands you a different feature for free.

\smallskip
\renewcommand{\arraystretch}{1.6}
\begin{tabular}{|l|c|l|}
\hline
\textbf{Form} & \textbf{Looks like} & \textbf{Shows you at a glance} \\
\hline
Standard & $ax^2 + bx + c$ & the $y$-intercept, which is \blank{1.4cm} \\
\hline
Vertex   & $a(x-h)^2 + k$  & the vertex, the point ( \blank{1cm} , \blank{1cm} ) \\
\hline
Factored & $a(x-r_1)(x-r_2)$ & the \blank{1.8cm} : $r_1$ and $r_2$ \\
\hline
\end{tabular}
\smallskip

From standard form, the \textbf{axis of symmetry} is the vertical line
$x = -\dfrac{b}{2a}$, and the vertex sits on it.

\smallskip
\textbf{Example 1.} Let $f(x) = x^2 - 6x + 5$. Find the vertex.

\begin{work}
  x &= -\frac{b}{2a} \\
    &= -\frac{-6}{2(1)} \\
    &= 3 \\
  f(3) &= (3)^2 - 6(3) + 5 \\
       &= -4
\end{work}

\begin{itemize}[itemsep=4pt]
  \item Vertex: ( \blank{1cm} , \blank{1cm} ) \qquad Axis of symmetry:
        $x = $ \blank{1.2cm}
  \item Vertex form: $f(x) = $ \blank{4cm}
  \item Factored form: $f(x) = (x-1)(x-5)$, so the zeros are \blank{1cm} and
        \blank{1cm}
  \item Here $a = 1 > 0$, so the parabola opens \blank{1.4cm} and the graph
        is concave \blank{1.4cm}.
\end{itemize}
\end{notesbox}

\vspace{0.1in}

\boxguard[30]
\begin{notesbox}{2. Average Rate of Change --- The Secant Line Returns}
From Lesson 1.3: the \textbf{average rate of change} of $f$ over $[a, b]$ is

\begin{center}
$\text{AROC over } [a,b] \;=\;
\dfrac{\text{change in \blank{1.6cm}}}{\text{change in \blank{1.6cm}}}
\;=\; \dfrac{f(b) - f(a)}{b - a}$
\end{center}

and it is the \blank{1.6cm} of the \textbf{secant line} through $(a, f(a))$
and $(b, f(b))$.

\smallskip
\textbf{Example 2.} For the same $f(x) = x^2 - 6x + 5$, find the AROC over
$[1, 4]$.

\begin{work}
  f(1) &= (1)^2 - 6(1) + 5 \\
       &= 0 \\
  f(4) &= (4)^2 - 6(4) + 5 \\
       &= -3 \\
  \text{AROC} &= \frac{-3 - 0}{4 - 1} \\
              &= -1
\end{work}

The dashed line below is that secant. Its slope is the AROC.

\begin{center}
\begin{tikzpicture}[x=0.72cm, y=0.32cm]
  \draw[linegray, very thin, step=1] (0,-5) grid (6,5);
  \draw[->, thick] (0,0) -- (6.6,0) node[below right] {\small $x$};
  \draw[->, thick] (0,-5) -- (0,5.6) node[above] {\small $f(x)$};
  \foreach \x in {1,2,3,4,5,6} \node[below] at (\x,0) {\small \x};
  \foreach \y in {-4,-2,2,4} \node[left] at (0,\y) {\small \y};
  \draw[very thick, plum] plot [smooth] coordinates
    {(0,5) (1,0) (2,-3) (3,-4) (4,-3) (5,0) (6,5)};
  \draw[dashed, thick, goldacc] (1,0) -- (4,-3);
  \fill[plum] (1,0) circle (2.5pt);
  \fill[plum] (3,-4) circle (2.5pt);
  \fill[plum] (4,-3) circle (2.5pt);
\end{tikzpicture}
\end{center}

\begin{itemize}[itemsep=4pt]
  \item AROC over $[3, 6]$: \quad $\dfrac{5 - (-4)}{6 - 3} = $ \blank{1.4cm}
  \item Same function, different interval, \blank{2.6cm} rate. A line could
        never do that --- and that difference is the rest of the lesson.
\end{itemize}
\end{notesbox}

\vspace{0.1in}

\boxguard[30]
\begin{notesbox}{3. Difference Tables --- Linear vs.\ Quadratic}
When the inputs are equally spaced, a \textbf{first difference} is one output
minus the one before it, and a \textbf{second difference} is one first
difference minus the one before it. With a step of $1$, each first difference
\textit{is} an AROC.

\medskip
\textbf{Linear:} $L(x) = 3x - 2$

\smallskip
\renewcommand{\arraystretch}{1.5}
\begin{tabular}{|c|c|c|c|c|c|}
\hline
$x$ & 0 & 1 & 2 & 3 & 4 \\
\hline
$L(x)$ & \blank{1cm} & \blank{1cm} & \blank{1cm} & \blank{1cm} & \blank{1cm} \\
\hline
\end{tabular}
\smallskip

First differences: \blank{1cm} , \blank{1cm} , \blank{1cm} , \blank{1cm}

Second differences: \blank{1cm} , \blank{1cm} , \blank{1cm}

For a \textbf{linear} function the AROC is \blank{2.2cm} over every interval,
so the second differences are all \blank{1cm}.

\medskip
\textbf{Quadratic:} $Q(x) = 2x^2 - 4x + 1$

\smallskip
\begin{tabular}{|c|c|c|c|c|c|}
\hline
$x$ & 0 & 1 & 2 & 3 & 4 \\
\hline
$Q(x)$ & 1 & $-1$ & 1 & 7 & 17 \\
\hline
\end{tabular}
\smallskip

First differences: \blank{1cm} , \blank{1cm} , \blank{1cm} , \blank{1cm}

Second differences: \blank{1cm} , \blank{1cm} , \blank{1cm}

The first differences are \textbf{not} constant --- but look at how they
change: each one is \blank{1cm} more than the one before. The AROCs of a
quadratic form a \blank{2cm} pattern of their own.

\medskip
\textbf{The headline:} for a quadratic, the average rates of change over
consecutive equal-length intervals change at a \blank{2.4cm} rate.
\end{notesbox}

\vspace{0.1in}

\boxguard[30]
\begin{notesbox}{4. Second Differences, the Leading Coefficient, and Concavity}
\textbf{The fingerprint.} If the inputs step by $h$, then a quadratic
$f(x) = ax^2 + bx + c$ has constant second differences equal to $2ah^2$.
When the step is $h = 1$, that constant is simply $2a$.

\begin{itemize}[itemsep=4pt]
  \item Check it on $Q(x) = 2x^2 - 4x + 1$: the second difference was
        \blank{1cm} , and $2a = 2(2) = $ \blank{1cm} --- they match.
  \item So a table alone gives you $a$: if the second differences are a
        constant $10$ with unit steps, then $a = $ \blank{1.2cm} .
\end{itemize}

\textbf{Concavity, from the numbers alone.} When the AROCs (first
differences) are \textit{increasing}, the graph is concave \blank{1.2cm};
when they are \textit{decreasing}, the graph is concave \blank{1.4cm}.

\smallskip
\renewcommand{\arraystretch}{1.6}
\begin{tabular}{|l|l|l|}
\hline
\textbf{Function} & \textbf{First differences} & \textbf{Second differences} \\
\hline
Linear    & constant                        & all \blank{1.2cm} \\
\hline
Quadratic & change by a \blank{1.6cm} step  & constant and \blank{1.8cm} \\
\hline
\end{tabular}
\smallskip

Since $Q$'s first differences $-2, 2, 6, 10$ are increasing, $Q$ is concave
\blank{1.2cm} --- and $a = 2$ is \blank{1.6cm} , which agrees.
\end{notesbox}

\vspace{0.1in}

\boxguard[24]
\begin{practicebox}
\begin{enumerate}[itemsep=8pt]
  \item Let $p(x) = -x^2 + 6x$. Complete the table and the difference rows.

        \smallskip
        \renewcommand{\arraystretch}{1.5}
        \begin{tabular}{|c|c|c|c|c|c|}
        \hline
        $x$ & 0 & 1 & 2 & 3 & 4 \\
        \hline
        $p(x)$ & \blank{1cm} & \blank{1cm} & \blank{1cm} & \blank{1cm} & \blank{1cm} \\
        \hline
        \end{tabular}
        \smallskip

        First differences: \blank{1cm} , \blank{1cm} , \blank{1cm} ,
        \blank{1cm}

        Second differences: \blank{1cm} , \blank{1cm} , \blank{1cm}

        The graph is concave \blank{1.4cm} , and $a = $ \blank{1.2cm} .

  \item A quadratic's AROCs over $[0,1]$, $[1,2]$, and $[2,3]$ are
        $1,\ 7,\ 13$.

        \begin{work}
          \text{step} &= 7 - 1 \\
                      &= 6 \\
          2a &= 6 \\
           a &= 3
        \end{work}

        \begin{enumerate}[label=(\alph*), itemsep=4pt]
          \item The AROCs step by \blank{1.2cm} each time, so the leading
                coefficient is $a = $ \blank{1.2cm} .
          \item The AROCs are increasing, so the graph is concave
                \blank{1.4cm} .
        \end{enumerate}
\end{enumerate}
\end{practicebox}

\end{document}
